\documentclass[11pt,a4paper]{article}
\usepackage[T1]{fontenc}
\usepackage[utf8]{inputenc}
\usepackage[brazilian]{babel}
\usepackage{amsmath,amssymb,booktabs,geometry,enumitem,listings,hyperref}
\geometry{margin=2.5cm}
\hypersetup{colorlinks=true,linkcolor=blue,urlcolor=blue}
\setlist{itemsep=3pt,topsep=4pt}
\lstset{basicstyle=\ttfamily\small,breaklines=true,columns=fullflexible,keepspaces=true,frame=single,showstringspaces=false}
\title{Guia de Estudo: Implementação do PAPT em Julia com JuMP e CPLEX}
\author{}
\date{}
\begin{document}
\maketitle
\tableofcontents
\newpage

\section{Objetivo deste guia}
Este material explica como partir da formulação matemática do Problema de Alocação de Professores a Turmas (PAPT) e construir uma implementação em Julia com JuMP e CPLEX. O código enviado pela orientadora serve como referência de organização, mas aparenta ser um trecho de um projeto maior, com conjuntos e parâmetros definidos fora do arquivo apresentado. A proposta é aproveitar a estrutura geral sem copiar diretamente as variáveis específicas daquele problema.

\section{Visão geral do código da orientadora}
O programa pode ser compreendido em etapas:
\begin{enumerate}
\item importar bibliotecas;
\item ler e organizar os dados da instância;
\item definir a função objetivo;
\item criar o modelo e as variáveis;
\item configurar o solver;
\item adicionar as restrições;
\item resolver o modelo;
\item extrair e registrar a solução.
\end{enumerate}

\section{Bibliotecas e parâmetros gerais}
O código começa com:
\begin{lstlisting}[language=Julia]
using JuMP, CPLEX, DataStructures, DelimitedFiles,
      Statistics, MathOptInterface, Combinatorics
\end{lstlisting}
\begin{center}
\begin{tabular}{p{0.30\textwidth}p{0.60\textwidth}}
\toprule
\textbf{Biblioteca} & \textbf{Função}\\
\midrule
\texttt{JuMP} & Escrever modelos matemáticos de otimização em Julia.\\
\texttt{CPLEX} & Interface para utilizar o solver IBM CPLEX.\\
\texttt{DelimitedFiles} & Ler arquivos de texto delimitados.\\
\texttt{MathOptInterface} & Interface padronizada entre JuMP e solvers.\\
\texttt{DataStructures}, \texttt{Statistics}, \texttt{Combinatorics} & Estruturas e operações auxiliares usadas em outras partes do projeto.\\
\bottomrule
\end{tabular}
\end{center}
O código também define:
\begin{lstlisting}[language=Julia]
const ZERO = 0.0001
const NTHREADS = 8
\end{lstlisting}
No trecho enviado, \texttt{ZERO} é usado como tolerância de gap relativo do CPLEX e \texttt{NTHREADS} define o número de threads. O gap é um critério de parada para o problema inteiro misto; não representa uma tolerância geral aplicada a todas as variáveis e restrições.
Para o primeiro protótipo do PAPT, comece com uma configuração simples e só depois experimente parâmetros avançados do solver. Assim, erros de formulação podem ser separados dos efeitos de configuração.

\section{Leitura e preparação dos dados}
O exemplo usa:
\begin{lstlisting}[language=Julia]
instance = ARGS[1]
const ninstance = parse(Int, instance)

filepath = "/home/pcandida/punim0040/instances/coordenadasSP.txt"
populationDP = readdlm(string(filepath, instance, "DP.txt"))
\end{lstlisting}
\texttt{ARGS[1]} recebe o primeiro argumento passado ao executar o script. Isso permite selecionar uma instância sem alterar o código. O caminho absoluto, porém, depende do computador da orientadora.
No PAPT, organize:
\begin{itemize}
\item $P$: professores;
\item $T$: turmas;
\item $K$: disciplinas;
\item $S$: slots semanais fixos;
\item $R$: combinações de dia e turno;
\item $h_j$: carga horária semanal da turma $j$;
\item $w_{ij}$: preferência do professor $i$ pela turma $j$;
\item $e_{ij}$: elegibilidade do professor $i$ para a turma $j$.
\end{itemize}
Separe dados institucionais observados dos parâmetros gerados para os experimentos. O professor registrado na oferta original é uma referência da alocação atual, não necessariamente a lista de todos os professores elegíveis para a turma.

\section{Função objetivo}
No código da orientadora, a função objetivo é encapsulada em:
\begin{lstlisting}[language=Julia]
function setObjective(m)
    # ...
    @objective(m, Min, ...)
end
\end{lstlisting}
A macro \texttt{@objective} define o sentido da otimização e a expressão a otimizar. O exemplo minimiza penalidades e custos. Para o PAPT, a versão não lexicográfica e sem normalização discutida anteriormente utiliza:
\begin{equation}
\max Z =
0{,}60F(x)+0{,}20\underline{s}
-0{,}10C(z)-0{,}10\overline{g}.
\end{equation}
Os critérios são:
\begin{align}
F(x) &= \sum_{i\in P}\sum_{j\in T} w_{ij}x_{ij},\\
C(z) &= \sum_{i\in P}\sum_{r\in R}z_{ir}.
\end{align}
Em JuMP:
\begin{lstlisting}[language=Julia]
@expression(model, F,
    sum(w[(i,j)] * x[i,j] for i in P for j in T)
)
@expression(model, C,
    sum(z[i,r] for i in P for r in R)
)
@objective(model, Max,
    0.60 * F +
    0.20 * s -
    0.10 * C -
    0.10 * g_max
)
\end{lstlisting}
Definir os critérios como expressões separadas facilita registrar seus valores individualmente. Como não há normalização, os pesos precisam ser avaliados, pois os critérios podem ter escalas diferentes.

\section{Variáveis de decisão}
No código recebido, as variáveis são criadas com \texttt{@variable}. A palavra \texttt{Bin} indica uma variável binária e \texttt{>= 0} define um limite inferior.
Para o PAPT:
\begin{lstlisting}[language=Julia]
@variable(model, x[i in P, j in T], Bin)
@variable(model, z[i in P, r in R], Bin)
@variable(model, g[i in P, k in K], Bin)
@variable(model, s >= 0)
@variable(model, g_max >= 0, Int)
\end{lstlisting}
\begin{center}
\begin{tabular}{p{0.20\textwidth}p{0.72\textwidth}}
\toprule
\textbf{Variável} & \textbf{Significado}\\
\midrule
$x_{ij}$ & Vale 1 se o professor $i$ for alocado à turma $j$; caso contrário, 0.\\
$z_{ir}$ & Vale 1 se o professor $i$ tiver aula em pelo menos uma turma no turno semanal $r$.\\
$g_{ik}$ & Vale 1 se o professor $i$ lecionar ao menos uma turma da disciplina $k$.\\
$\underline{s}$ & Piso da satisfação dos professores incluídos no critério de equidade.\\
$\overline{g}$ & Maior número de disciplinas distintas atribuído a um professor.\\
\bottomrule
\end{tabular}
\end{center}
Os nomes \texttt{s} e \texttt{g\_max} são nomes de implementação e representam $\underline{s}$ e $\overline{g}$.

\section{Configuração do solver}
O código recebido configura o CPLEX com:
\begin{lstlisting}[language=Julia]
set_optimizer_attribute(m, "CPX_PARAM_THREADS", NTHREADS)
set_optimizer_attribute(m, "CPX_PARAM_EPGAP", ZERO)
set_optimizer_attribute(m, "CPX_PARAM_TILIM", 3600)
\end{lstlisting}
Esses parâmetros controlam threads, gap relativo e limite de tempo. As opções de heurísticas e cortes são configurações do solver, não restrições matemáticas.
Para o primeiro protótipo, mantenha configurações essenciais e introduza ajustes avançados somente depois de validar o modelo.

\section{Tradução das restrições para JuMP}
A macro \texttt{@constraint} permite declarar restrições. A cobertura exige:
\begin{equation}
\sum_{i\in P}x_{ij}=1,\qquad \forall j\in T.
\end{equation}
Em JuMP:
\begin{lstlisting}[language=Julia]
@constraint(model, cobertura[j in T],
    sum(x[i,j] for i in P) == 1
)
\end{lstlisting}
A expressão \texttt{j in T} cria uma restrição para cada turma.
\begin{center}
\begin{tabular}{p{0.34\textwidth}p{0.56\textwidth}}
\toprule
\textbf{Restrição matemática} & \textbf{Tradução conceitual}\\
\midrule
Cobertura & Exatamente um professor por turma.\\
Elegibilidade & $x_{ij}\leq e_{ij}$.\\
Carga mínima & A soma das horas atribuídas ao professor é pelo menos $h_i^{\min}$.\\
Carga máxima & A soma das horas atribuídas ao professor não ultrapassa $h_i^{\max}$.\\
Piso de satisfação & $\underline{s}\leq\sum_j w_{ij}x_{ij}$ para cada professor considerado.\\
Ligação turma--disciplina & $x_{ij}\leq g_{i,\mathrm{Disc}(j)}$.\\
Teto de disciplinas & $\overline{g}\geq\sum_k g_{ik}$.\\
Conflito por slot & Cada professor assume no máximo uma turma em cada slot fixo.\\
Compactação por turno & Se uma turma alocada ocorre no turno $r$, então $z_{ir}=1$.\\
\bottomrule
\end{tabular}
\end{center}
As ligações reversas de $g_{ik}$ e $z_{ir}$ também devem ser implementadas. Elas garantem que as variáveis auxiliares representem a alocação real, em vez de permanecerem artificialmente maiores que a ocupação efetiva.

\section{Exportação, resolução e leitura da solução}
O código da orientadora exporta o modelo:
\begin{lstlisting}[language=Julia]
write_to_file(m, "./model/ps" * instance * ".lp")
\end{lstlisting}
O arquivo \texttt{.lp} permite inspecionar a formulação construída. A resolução é iniciada com:
\begin{lstlisting}[language=Julia]
optimize!(model)
\end{lstlisting}
Após a resolução, verifique o status do solver antes de acessar valores das variáveis. Para benchmarks, registre:
\begin{itemize}
\item status da otimização;
\item valor da função objetivo;
\item alocação de cada turma;
\item carga semanal de cada professor;
\item satisfação individual;
\item número de turnos ocupados;
\item máximo de disciplinas por professor;
\item tempo de execução e gap, quando disponíveis.
\end{itemize}

\section{Observações sobre o arquivo recebido}
O trecho enviado aparenta depender de código que não está incluído. Conjuntos e parâmetros como \texttt{K}, \texttt{T}, \texttt{P}, \texttt{priority} e \texttt{demand} são utilizados sem definição no fragmento. Além disso, alguns acessos às variáveis \texttt{V} e \texttt{Q} aparentam usar quantidades de índices diferentes das declarações apresentadas. Esses pontos devem ser conferidos no projeto completo. Assim, o arquivo é mais útil como referência de organização do que como programa autônomo.

\section{Plano de desenvolvimento recomendado}
\begin{enumerate}[label=\textbf{Etapa \Alph*.}]
\item \textbf{Instância pequena fixa.} Codificar diretamente em Julia os três professores e as quatro turmas da instância reduzida.
\item \textbf{Validação matemática.} Conferir cobertura, elegibilidade, carga horária, conflitos, satisfação e variáveis auxiliares.
\item \textbf{Entrada JSON.} Depois que o protótipo funcionar, carregar o JSON normalizado produzido pelo Python.
\item \textbf{Benchmarks.} Automatizar leitura das instâncias, construção do modelo, resolução e registro dos resultados.
\end{enumerate}
Essa sequência reduz fontes possíveis de erro. Com dados escritos diretamente no código, divergências tendem a estar na tradução matemática ou nas restrições. A leitura do JSON pode ser adicionada depois, como etapa independente.

\section{Primeiro teste e valores de referência}
O exemplo inicial utiliza a instância reduzida e a função objetivo ponderada sem normalização. Os valores de referência para a solução indicada no estudo anterior são:
\begin{center}
\begin{tabular}{lc}
\toprule
\textbf{Indicador} & \textbf{Valor de referência}\\
\midrule
Preferência total $F$ & 17\\
Piso de satisfação $\underline{s}$ & 3\\
Turnos ocupados $C$ & 8\\
Máximo de disciplinas $\overline{g}$ & 1\\
Função objetivo $Z$ & 9{,}9\\
\bottomrule
\end{tabular}
\end{center}
Esses valores servem como referência para conferir a implementação; não constituem confirmação de que o arquivo foi executado nesta etapa. A execução depende de Julia, JuMP, CPLEX e de uma instalação/licença compatível do solver.

\section{Checklist de estudo}
\begin{itemize}[label=$\square$]
\item Sei explicar a função de JuMP e CPLEX.
\item Sei distinguir conjuntos, parâmetros e variáveis de decisão.
\item Consigo traduzir a restrição de cobertura para \texttt{@constraint}.
\item Consigo explicar por que $x_{ij}$, $z_{ir}$ e $g_{ik}$ são diferentes.
\item Sei por que o modelo deve ser validado com uma instância pequena.
\item Sei quais indicadores registrar em cada benchmark.
\end{itemize}
\end{document}
